\storyheader{Người chiến thắng}{Ngày 01}{

Tại một ngôi trường hết sức bình thường tên là trường cao trung Seiran ở
một đất nước hết sức bình thường là Nhật Bản với những học sinh cũng hết
sức bình thường đang theo học tại nơi đây. Tuy nhiên, ở nơi này lại diễn
một chuyện không được bình thường cho lắm, một tình tiết tưởng chừng như
chỉ diễn ra ở những bộ shoujo cổ điển.

Shin, Haruka và Akira là ba chàng trai năm cuối có thể gọi là nằm ở tầng
lớp ``thượng lưu'' ở đây. Học giỏi có, đẹp trai có và tốt bụng cũng có,
khiến cho gần như mọi nữ sinh trong trường từ bạn học đến những đàn em
lớp dưới vô cùng mến mộ. Và kể cả những nam sinh cũng thầm nể phục ba
người này. Thế nhưng định mệnh trớ trêu thay lại khiến cho họ cùng nhau
phải lòng một cô bạn cùng lớp tên là Mei.

Mei là một cô gái vô cùng điển hình của những nữ chính shoujo manga mang
motip harem ngược. Gương mặt gần như là vô cảm, đôi mắt đen láy ánh lên
sự lạnh lùng và thờ ơ với tất cả mọi người, mái tóc dài đen tuyền được
để xõa phía sau lại càng tô điểm thêm cho vẻ đẹp phi giới tính của cô.
Và đúng như ngoại hình, Mei là một người lạnh lùng và rất ít nói. Không
phải cô nhút nhát, mà cô chỉ đơn giản là không hề quan tâm đến thứ gì
xung quanh mà thôi. Và chính tính cách ấy đã hoàn toàn mê hoặc ba chàng
trai kia, khiến cho họ nảy sinh một thứ tình cảm lãng mạn với cô và
quyết tâm làm tan chảy trái tim băng giá ấy.

Thế nhưng như đã nói ở trên, Mei hoàn toàn không thèm để tâm đến bất cứ
thứ gì ngoài cuốn sách luôn cầm khư khư trên tay cả, thế nên sau một
thời gian thì cả ba đều nhận ra mọi nỗ lực của họ hoàn toàn là vô nghĩa.

Tuy nhiên lúc này bọn họ lại đồng loạt nghĩ đến cùng một người có thể sẽ
trở thành cứu tinh của họ, cô bạn hoạt bát vô cùng nổi bật trong lớp
mang tên Hikaru, và cũng là bạn thuở nhỏ của Mei.

Nếu ngoại hình của Mei mang một nét đẹp phi giới tính, thì vẻ đẹp của
Hikaru có thể gọi là phi thực với mái tóc trắng ngần lai giữa Nhật và
Nga được để xõa như thể một dòng suối tuyết trắng xóa, đôi mắt màu xanh
ngọc bích long lanh như thể một hồ nước tuyệt đẹp, và làn da trắng nõn
nà lại càng khiến cho Hikaru dường như không thuộc thế giới này. Và tất
nhiên với ngoại hình như vậy thì Hikaru cũng là một cô gái rất nổi bật
trong trường, nhưng trái ngược với Mei thì Hikaru lại là kiểu người rất
hoạt bát và có thể xem là nguồn năng lượng của tất cả mọi người.

Tuy nhiên, điểm mà mọi người nể phục nhất ở Hikaru đó là việc cô có thể
dễ dàng xâm nhập vào cái thế giới nội tâm gần như là tách biệt của Mei,
và cô cũng là người duy nhất có thể khiến cho Mei chịu bỏ cuốn sách
xuống và nói chuyện một cách đàng hoàng, đồng thời bày ra những biểu cảm
ấm áp dường như là chỉ dành riêng cho Hikaru.

Vì lý do đó nên cả Shin, Haruka và Akira đã cùng đi đến một kết luận
chung mà thậm chí còn chả cần bàn bạc, gọi riêng Hikaru ra thư viện để
nhờ cô giúp đỡ trong việc làm tan chảy trái tim băng giá của Mei.

Liệu ai trong số họ sẽ giành được trái tim của Mei?

Hồi sau sẽ rõ!

*

Lại còn hồi sau sẽ rõ luôn cơ, kịch bản với giọng văn đúng bài mấy bộ
manga shoujo hồi xưa luôn ấy chứ.

Mà chắc cả đời tôi cũng không thể nào ngờ được bản thân sẽ có thể sống
được đến cái ngày mà mấy cái nội dung rồi tình tiết sến súa đó được đem
ra ngoài thực tế luôn quá, nhưng rốt cuộc thì tôi cũng đã sống được đến
thời khắc đó rồi, cụ thể là qua ba cậu bạn quen thuộc ở trước mặt đây.

Và tôi thì không ai khác chính là cô bạn xinh đẹp tài năng, luôn mang
dáng vẻ tràn đầy năng lượng và là nguồn chữa lành cho tất cả mọi người,
mang một nét đẹp đến phi thực tế, người sẽ giúp ba nam chính tiến tới
với cô bạn lạnh lùng của mình, Hikaru!

...Thôi, vậy được rồi, nghiêm túc thôi nào.

``Thế tóm lại là mấy cậu đang bế tắc trong việc lấy lòng Mei, nên là nhờ
mình giúp đỡ ấy hả?'' Tôi ngồi lên ghế, lấy tay chống cằm và hỏi ba cậu
bạn đứng trước mặt.

``Ừm thì, như cậu vừa nghe thấy đấy, mình muốn có sự chú ý của Mei nhưng
hoàn toàn bế tắc, và rồi tìm đến cậu.'' Shin lên tiếng, đồng thời cậu ta
cũng đánh mắt liếc sang hai người còn lại. ``Còn về hai người này thì...
chắc cũng giống vậy.''

``Chẳng thể ngờ ba người chúng ta lại đi thích cùng một người ha?''
Haruka cười phá lên với cái điệu bộ phô trương vốn có của mình. ``Mà, dù
sao thì người thắng cũng sẽ là tôi thôi!''

``Bớt điêu đi, nếu đã phải vác xác đến đây làm phiền Hikaru thì mày cũng
đang bí vãi ra thôi.'' Akira khoanh tay và thẳng thừng phủ định cái
tương lai do Haruka vẽ ra.

Đúng như trong cái cốt truyện vừa được xóa bỏ khỏi đầu tôi vài giây
trước thì ba người này chính là ba nam sinh phải gọi là nổi như cồn ở
trường và cùng thích một người là cô bạn thuở nhỏ lạnh lùng tên là Mei
của tôi. Shin là hội trưởng câu lạc bộ văn học với tính cách khá thân
thiện và lịch sự, đội trưởng đội điền kinh Haruka mang tính cách rất chi
là khoa trương và ngạo mạn.

Hai con người kia thì là tuýp nhân vật khá quen thuộc rồi, nhưng người
khiến tôi bất ngờ là cậu bạn Akira. Cậu ta là một thành viên rất quen
thuộc của câu lạc bộ về nhà, giọng điệu thì thô lỗ, tính cách thì bất
cần, tuy nhiên về cơ bản vẫn khá là tốt tính. Tôi không nghĩ một người
như cậu ta lại quan tâm đến mấy chuyện như này đấy, đúng là tình yêu sẽ
làm con người ta thay đổi mà.

``Tất nhiên thì cậu cũng sẽ không thiệt. Nếu cậu cần giúp đỡ việc gì thì
chúng mình sẽ luôn sẵn lòng.''

``Thật ra cũng không cần phải làm đến vậy đâu, mình luôn sẵn sàng giúp
đỡ mà, Mei có thêm bạn là chuyện tốt chứ sao.''

Về cơ bản đã học chung với nhau tận ba năm trời của thời cao trung, cũng
trải qua đủ sự kiện cùng nhau rồi, thế nên tôi biết rõ rằng họ đều là
người tốt hết cả và sẽ luôn sẵn lòng giúp đỡ nếu ba người này thật lòng
muốn làm bạn với Mei.

``Tuy nhiên...''

Nếu ba người họ đã muốn làm người yêu thay vì bạn bè của Mei, thì tôi
cũng không thể cứ thế mà giúp họ không lý do được. Tôi muốn Mei được
hạnh phúc, và sẽ tuyệt đối không chấp nhận giúp đỡ những người nào có
tình cảm hời hợt.

``Các cậu có thể nói cho mình biết là các cậu thích Mei ở điểm nào
không?'' Tôi hỏi ba chàng trai ở trước mặt.

Trong những kịch bản của shoujo manga, các nhân vật nam chính ban đầu
chỉ xem cô nữ chính kiêu kỳ là những ``viên ngọc quý'', những ``bông hoa
kỳ lạ'' khiến cho bọn họ nổi lên mong muốn chinh phục và dần sẽ phát
triển thành tình yêu thực sự. Tuy nhiên dù thế nào đi nữa thì tôi cũng
không thích thú gì với việc Mei bị xem như một thứ để chinh phục đâu.

Nào, các cậu sẽ trả lời sao đây?

``Hỏi gì hiển nhiên vậy? Tất nhiên là vì tính cách của cậu ấy rồi.''
Akira nhanh chóng lên tiếng với giọng điệu khá là ngạc nhiên như thể
không hiểu sao tôi lại hỏi câu này. ``Một người sống thật với bản thân,
không hề quan tâm người khác nghĩ gì về mình mà chỉ sống theo cách bản
thân mong muốn, chẳng phải như vậy là quá dư lý do rồi sao?''

``Ai cũng sẽ thích một cô nàng mạnh mẽ và sống thật với bản thân như vậy
mà ha?'' Haruka lên tiếng ngay sau đó để củng cố cho sự khẳng định chắc
nịch của Akira.

``Bị tranh mất lời nói như này quả là khó chịu thật đấy.'' Shin gãi đầu
và cười gượng ``Nhưng mà ừ, mình cũng giống hai người họ thôi, muốn ở
bên và đồng hành với một người mạnh mẽ như vậy thì đâu có gì sai nhỉ?''

Vậy là, họ bị thu hút bởi dáng vẻ mạnh mẽ, luôn sống thật với bản thân
của Mei nhỉ?

Nghe khá tương đồng với...

Mà thôi bỏ đi.

Dù sao thì có lẽ ba người họ cũng không quá tệ, với lý do đó thì tôi
hoàn toàn có thể hiểu và đồng cảm được.

``Thôi được rồi, mình sẽ giúp các cậu một tay, cùng cố gắng nhé!''

Tôi lên tiếng đồng ý với lời đề nghị ban đầu của họ. Mặc kệ cho con tim
vẫn đang gào lên phản đối.

*

Được rồi, phải thú nhận thôi, là giờ tôi đang thực sự cảm thấy vô cùng
là hối hận với quyết định của bản thân hồi trước.

Đã một thời gian trôi qua rồi kể từ khi tôi chấp nhận giúp đỡ ba người
kia, và thật sự đó là khoảng thời gian khiến tôi cảm thấy mệt mỏi nhất
trên đời. Trong suốt khoảng thời gian đó, tôi đã liên tục đưa ra những
lời khuyên, chỉ cho họ biết sở thích của Mei là như thế nào, đôi khi
cũng sắp xếp cho một trong ba người họ những khoảnh khắc được ở riêng
với Mei. Và thật sự nó khá hiệu quả khi bây giờ Mei đã có nhiều bạn hơn,
đúng với ý định của tôi ban đầu.

Tuy nhiên, sai lầm của tôi là việc nghĩ cả ba người họ sẽ có thể xoay sở
ngon ơ khi được tôi mớm cho nhiều thông tin quan trọng. Thế nhưng cả đời
tôi không bao giờ có thể tưởng tượng được rằng ba chàng trai hoàn hảo
của trường khi vướng vào chuyện tình cảm thì lại cực kỳ ngây thơ. Họ như
kiểu những quả bom nổ chậm luôn chực chờ để đồng lòng spam đầy hộp thư
điện thoại của tôi. Mấy tin nhắn như kiểu ``Hikaru ơi cứu mình với! Thơ
mình viết tặng Mei bị cô ấy ném ra sọt rác rồi!'', ``Ê Hikaru, tại sao
khi tôi biểu diễn kỹ thuật đá bóng thần sầu của mình thì Mei lại bỏ về
giữa chừng vậy?'' hay ``Hikaru, tôi nên tặng cho Mei cuốn sách vật lý
lượng tử hay triết học Mác Lê Nin đây?'' thường xuyên nằm đầy ú ụ trong
hộp tin nhắn của tôi.

Và tôi, một cô gái xinh đẹp với tấm lòng siêu tốt bụng tất nhiên sẽ giúp
đỡ họ hết sức mình, nhưng mà cái gì cũng vừa vừa phai phải thôi, tôi đây
thì cũng chỉ là con người thôi mà! Ban đầu thì tôi còn nhiệt tình tư
vấn, nhưng một thời gian thì não của tôi thật sự đã đạt đến giới hạn
chịu đựng rồi. Với cả còn chưa kể đến cái tôi cao ngút trời của ba tên
kia đâu nhé, tên nào tên nấy cũng háo thắng y như nhau, khiến cho tôi
lắm lúc còn tưởng mình đang là người hỗ trợ cho một cái chương trình
truyền hình thực tế dở tệ nào đó không đấy!

``Hầy...'' Tôi nằm dài ra bàn với vẻ mệt mỏi và uể oải.

Đột nhiên tôi cảm thấy có một ngón tay của ai đó chọc chọc vào đầu mình
nên là đã ngẩng dậy xem. Trước mặt tôi là gương mặt vô cảm mang nét đẹp
phi giới tính quen thuộc mà tôi luôn nhìn thấy mỗi ngày giờ đây mang một
chút vẻ lo lắng hiện rõ trên ánh mắt.

Đây chính là Mei, cô bạn thuở nhỏ của tôi mà ba người kia đang đánh nhau
sứt đầu mẻ trán để mà có được.

``Sao vậy?''

``À, mình chỉ hơi mệt chút thôi, cậu đừng lo quá...'' Tôi trả lời Mei
khi vẫn đang nằm dài ra trên bàn.

Mei nghe xong thì im lặng một hồi lâu, rồi cậu ấy không nói không rằng
đi ra ngoài mất, bỏ lại tôi với một sự hoang mang không hề nhỏ. Sau
khoảng 10 phút thì cậu ấy đã quay lại, trên tay là một lon nước cam.

``Đây, cho cậu.'' Mei đặt lon nước lên bàn trước mặt tôi.

``Cảm ơn nhé!''

Mei thường không hay thể hiện những cử chỉ thân thiện như này với ai
khác mà chủ yếu là sẽ bỏ đi luôn và không quan tâm. Thế nên tôi có thể
xem là nằm ở vị trí độc tôn và không thể thay thế trong lòng của Mei.
Tôi cũng rất tự hào về việc này đấy nhé.

Thế nhưng quay lại vấn đề chính, thật sự thì tôi không thể nào mà cùng
một lúc chỉ đạo cả ba người đối đầu với nhau đến tận khi mà một trong
mấy người bọn họ được lựa chọn bởi Mei được. Tôi nghĩ rằng mình bắt buộc
phải có một cách nào đó để đẩy nhanh tiến độ lên mới được, và cũng phải
đưa ra giới hạn thời gian cho cái sự giúp đỡ này thôi, não tôi thật sự
đã đến giới hạn rồi. Vì vậy nên chiều hôm đó khi tan học, tôi kéo cả Mei
lẫn ba người kia ra phía sau trường.

``Sao tự nhiên lại kéo cả bọn ra đây vậy? Có việc gì sao?''

Tôi không trả lời câu hỏi của Akira, mà thay vào đó lại quay sang Mei.

``Mei, ba người này đang thích cậu đấy.''

``Ê này!!!'' Cả ba người đều phát hoảng lên khi thấy tôi thản nhiên tiết
lộ như vậy. ``Sao lại nói ra cho cậu ấy biết thế hả!?''

``Mình muốn xác nhận xem thật sự Mei có thích ai trong mấy cậu chưa,
chúng ta đã gần tốt nghiệp rồi còn gì? Trong tận một năm qua nếu mà Mei
vẫn không thích ai thì thật sự bó tay chịu trận rồi đấy.''

``Cũng đúng, nhưng mà...''

``Vậy Mei, cậu có đang thích ai không?'' Tôi quay sang hỏi cô bạn bên
cạnh mình.

Hỏi thế thôi chứ thực sự là tôi không nghĩ cậu ấy sẽ đổ một trong ba
người kia đâu, vì về cơ bản là tôi chỉ giúp mấy cậu ấy bằng cách nói cho
họ biết sở thích của Mei là gì, cách nói chuyện của Mei và cho họ vài
khoảng thời gian riêng tư với Mei chứ còn lại vẫn là tự lực gánh sinh
thôi, và như mọi người biết rồi đấy, họ luôn làm hỏng cơ hội mà tôi trao
cho và nhận thất bại đau điếng. Thế nên trong số ba người họ thì Mei
chắc sẽ không thích ai đâu.

``Mình có thích một người.''

...Ể?

``Và thật ra thì trong lòng của mình cũng đã có câu trả lời từ lâu
rồi.'' Mei tiếp tục nói. ``Mình sẽ công bố trong buổi lễ tốt nghiệp diễn
ra vào ngày mai.''

*

Giờ nghĩ lại, hồi xưa thì tôi đúng là mặt dày vô đối luôn. Khi học lớp
một thì tôi đã luôn chú ý đến một cô bạn kỳ lạ ở trong lớp. Cậu ấy không
thèm nói chuyện với ai, cũng không không muốn ai tiếp cận và nói chuyện
với mình, khiến bầu không khí xung quanh cô bạn ấy như thể có một bức
tường ngăn cách hoàn toàn với thế giới bên ngoài vậy.

Tuy nhiên khác với mọi nguời khi đó e ngại và tránh nói chuyện với Mei,
thì tôi lại bị hút hồn bởi sự mạnh mẽ đó, nên là cứ liên tục mặt dày trò
chuyện với Mei không biết bao nhiêu lần, cho đến tận khi mà tôi trở
thành người duy nhất thân thiết và bắt được nhịp nội tâm của cậu ấy.

Có lẽ tôi đã yêu Mei ngay từ cái nhìn đầu tiên, và khoảng thời gian ở
bên cậu ấy đã dần khuếch đại cảm xúc bên trong tôi ngày một lớn hơn cho
đến khi tôi hoàn toàn nhận thức được rằng tình cảm đó không phải là tình
bạn.

Sau một khoảng thời gian vừa đi vừa suy nghĩ vu vơ thì thoáng cái mà tôi
đã về được đến căn phòng quen thuộc của mình. Tôi đóng cửa lại, thở dài
một tiếng nhẹ và thả người nằm phịch xuống giường với trái tim vẫn đang
đau nhói của mình.

``Hầy, cứ tưởng mình là người thắng chứ nhỉ? Mấy cậu ấy coi bộ cũng ra
trò thật đấy.'' Tôi lẩm bẩm trong miệng. ``Mà, cũng là do mình thôi.''

Tuy là mang trong mình thứ tình cảm sai trái này với Mei thì tôi cũng
chưa bao giờ thổ lộ nó với cô bạn của mình. Một phần vì cả hai đều là
con gái, nhưng phần lớn là vì Mei.

Tôi biết rằng Mei chỉ đơn giản xem tôi là một chỗ dựa vững chắc, một
người bạn quan trọng hiếm hoi trong thế giới nội tâm đầy cô độc của cậu
ấy mà thôi. Và tôi rất hạnh phúc với điều đó, thực sự vô cùng hạnh phúc.
Vì việc được ở bên Mei, được nhìn thấy cậu ấy bỏ cuốn sách xuống để lắng
nghe tôi, được nhìn thấy nụ cười ấm áp chỉ dành riêng cho tôi thì cũng
đã đủ lắm rồi. Những khoảnh khắc ấy với tôi mà nói giống như liều thuốc
an thần, giúp bản thân có thể quên đi phần nào mối tình đơn phương của
mình.

Đối với bản thân tôi thì tình bạn với Mei là điều quý giá nhất giữa tôi
và cậu ấy. Và tôi sẵn sàng làm bất cứ chuyện gì để giữ gìn được tình bạn
ấy, dù điều đó đồng nghĩa với việc phải tự tay bóp nghẹt trái tim của
chính mình.

Thế nhưng rốt cuộc thì tôi cũng vẫn chỉ là một con người yếu đuối mà
thôi. Trong quá trình hỗ trợ Shin, Haruka và Akira, dù tôi đã cố gắng lờ
đi cảm xúc của bản thân thì tôi vẫn cảm giác có đôi chút chạnh lòng. Mỗi
lần thấy Mei khẽ mỉm cười với một món quà nào đấy từ ba người kia, hay
trả lời dài hơn một chút trong cuộc trò chuyện của họ là mỗi lần mà vết
thương đáng lẽ phải lành trên trái tim tôi lại tiếp tục rỉ máu.

Tôi biết chứ, tôi biết rõ đối với Mei thì tôi vẫn luôn là người bạn quan
trọng nhất trên đời, là người duy nhất khiến cho Mei thật sự cảm thấy
thoải mái khi ở bên, là người duy nhất mà Mei ôm chầm lấy một cách không
do dự. Tuy nhiên thì rồi sẽ có một ngày không xa nào đó, tôi sẽ đánh mất
vị trí ấy trong lòng Mei, và giờ thì cái ngày không xa đó sẽ đến chỉ
trong ngày mai mà thôi.

Trong những tựa game otome hay những bộ shoujo manga harem ngược mà
trước kia tôi từng mê mẩn thì nữ chính sẽ luôn luôn đến với các nam
chính. Với game thì sẽ chia ra rất nhiều tuyến đường để bạn lựa chọn xem
nên về với ai, còn với manga thì sẽ có những tác phẩm chia ra nhiều
ending khác nhau và nữ chính sẽ đến với tất cả nam chính. Tuy nhiên thì
``nhân vật bạn thân'' dù có tình cảm với nữ chính đi chăng nữa thì cũng
chỉ là một thứ phụ gia giúp đóng góp vào mạch câu chuyện thôi, chẳng có
một cái kết nào là dành cho họ cả.

``Haha, sao tự nhiên hôm nay mình đa cảm thế nhỉ?'' Tôi ngồi dậy và lắc
đầu, cố gắng xua tan đi cảm xúc kỳ lạ trong lòng. ``Mei xứng đáng được
hạnh phúc, và mình sẽ là người chứng kiến điều đó với nụ cười chân thành
nhất.''

Thế nhưng rốt cuộc thì lại chẳng đâu vào đâu cả, tôi không thể nào cân
bằng được cảm xúc bên trong mình nữa rồi.

Đúng là tệ thật mà...

Tôi vớ lấy con gấu bông được đặt ở góc phòng. Đây là món quà sinh nhật
đầu tiên mà Mei tặng cho tôi, thế nên cũng có thể nói rằng toàn bộ kỷ
niệm trong suốt hơn mười hai năm gắn bó giữa tôi và Mei đều nằm gói gọn
bên trong con gấu bông này hết.

Khi ôm lấy con gấu bông ấy, một giọt nước ấm nóng chảy ra từ khóe mắt,
lăn dài trên má tôi.

Ôm một mối tình đơn phương như vậy suốt hơn mười hai năm dài đằng đẵng
mà không kêu ca gì cả, tôi cũng xứng đáng được khen ngợi đấy chứ? Tôi
rất mạnh mẽ mà nhỉ?

Ngày mai, tôi sẽ lại trở thành Hikaru của mọi khi. Một cô gái hoạt bát,
thân thiện, luôn sẵn sàng giúp đỡ bạn bè và tiếp tục đóng vai trò là
nguồn năng lượng chữa lành cho tất cả mọi người trong buổi lễ tốt nghiệp
buồn bã và đầy nước mắt sắp tới.

Thế nên lúc này, chỉ giây phút này thôi, tôi có quyền để con tim của
mình yếu đuối một chút mà phải không?

*

Sau buổi lễ tốt nghiệp thấm đẫm nước mắt buồn bã của những cuộc ly biệt
và nước mắt hạnh phúc của những người vui mừng khi cuối cùng cũng thoát
khỏi ngôi trường này thì mọi người cũng đã tản ra, với những ước mơ
tương lai sẽ được viết tiếp trên trang giấy mang tên đại học.

Và cũng chính lúc này đây là lúc mà tôi đặt dấu chấm cho mối tình đơn
phương kéo dài hơn mười hai năm của mình.

Tôi lặng lẽ đi vòng ra sau trường. Tại đó có một cái cây hoa anh đào
đang nở rộ như thể muốn tiễn biệt một thế hệ học sinh nữa sắp ra khỏi
mái trường, và cũng là nơi Mei đang đứng để đợi ba người kia đến. Cố
kiềm lại mong muốn ngắm hoa anh đào, tôi ẩn mình sau một cái bụi cỏ gần
gốc cây, thân hình nhỏ bé của tôi dễ dàng bị nó che khuất.

Tôi cũng không biết rõ mục đích của bản thân khi ở đây là gì nữa, có thể
là để làm người đầu tiên chúc mừng Mei chăng? Để nhào tới ôm lấy cậu ấy
và nói rằng ``Chúc cậu hạnh phúc nhé Mei!''. Hoặc cũng có thể là để nhìn
thấy nụ cười của Mei dành cho một ai khác, cái nụ cười ấm áp hiếm hoi
trên gương mặt vô cảm mà đã từng chỉ dành riêng cho một mình tôi.

Gió thổi thoảng qua mang theo hương thơm nhàn nhạt của hoa anh đào, tôi
hít một hơi thật sâu và cố gắng giữ bình tĩnh. Sẽ ổn thôi, khi Mei đưa
ra câu trả lời, tôi cũng sẽ bước ra và chúc mừng cậu ấy một cách chân
thành nhất có thể.

Thế nhưng mười phút, ba mươi phút, và rồi tận một tiếng trôi qua, mặt
trời cũng đang dần lặn sau những tòa nhà cao tầng ở Tokyo, mang đến sắc
cam nhuộm toàn bầu trời, thế nhưng ba người kia vẫn chưa xuất hiện. Tôi
nhíu mày, một chút giận dữ dâng lên trong lòng.

``Sao họ lại dám để Mei đợi thế này?'' Tôi lẩm bẩm với giọng điệu bực
tức. ``Mình sẽ cho họ biết tay, dám làm Mei buồn!''

Mei vẫn đang đứng đợi một mình dưới gốc cây anh đào, khuôn mặt thanh tú
không để lộ một chút biểu cảm nào. Nhưng tôi biết rằng Mei luôn đúng giờ
và sự chậm trễ của người khác sẽ khiến cậu ấy khó chịu.

Lúc này, tôi chỉ muốn lao ra mắng ba tên kia một trận, muốn hét lên rằng
họ không xứng, rằng họ dám để Mei đợi thế này thì đừng mơ có được cậu ấy
nữa.

Nhưng tôi không có tư cách, tôi chỉ là người đứng sau bụi cỏ, nhìn người
mình yêu nhất thế giới đứng đợi người khác mà thôi.

Thời gian cứ thế mà trôi qua một cách vô cùng chậm chạp khiến lòng tôi
càng thêm sốt ruột. Lặng lẽ liếc nhìn Mei qua kẽ lá, tôi thấy cậu ấy vẫn
đang yên lặng chờ đợi, tay mân mê một cánh hoa anh đào vừa rơi xuống
khiến trái tim tôi lại một lần nữa nhói đau, khiến tôi lúc này chỉ muốn
lao ra đó và ôm lấy Mei, muốn nói rằng ``Đừng đợi nữa, về với mình đi!
Mình sẽ không để cậu phải đứng một mình thế này đâu!''

Thế nhưng, đúng lúc đó, Mei cũng đã bắt đầu lên tiếng.

``Mình biết cậu ở đó, Hikaru.''

Tim tôi hẫng đi một nhịp khi nghe câu nói ấy.

Cậu ấy biết.

Từ bao giờ?

Từ lúc tôi vừa núp vào bụi cỏ?

Hay từ lúc tôi còn đang đứng ngẩn ngơ giữa sân trường, đấu tranh xem có
nên đến đây không?

Tôi cứng đờ người, không dám thở mạnh. Cánh hoa anh đào vẫn rơi, gió vẫn
thổi, nhưng mọi thứ xung quanh như bị bấm nút pause.

Mei không quay lại, cậu ấy vẫn đứng đó, lưng thẳng, mái tóc dài bị gió
thổi bay, giọng nói tiếp tục vang lên trong không gian chỉ còn lại hai
chúng tôi.

``Cậu muốn ra hay không, tùy cậu. Nhưng chí ít thì xin cậu hãy nghe hết,
làm ơn đừng bỏ chạy nhé.''

*

Tôi không biết bản thân đã phải lòng Hikaru từ khi nào.

Hồi đó, những từ ngữ như trầm tính, ít nói, khó gần, đáng sợ là những từ
ngữ mà mọi người thường dùng để miêu tả tôi. Ở trường mẫu giáo, bạn bè
chẳng ai chịu chơi chung với tôi cả, họ tránh né tôi như thể tránh một
đám mây đen, thì thầm sau lưng rằng không khí xung quanh tôi lúc nào
cũng khó gần và chẳng có một chút ấm áp nào cả. ``Mei chẳng bao giờ mỉm
cười cả, nhìn mặt của cậu ấy thôi là đã thấy đáng sợ rồi.'' Một cậu bạn
ngày xưa đã từng nói như thế, và nó như thể một vết dao cứa vào trái tim
non nớt của tôi ngày đó vậy.

Khi đó, tôi không thể hiểu được tại sao mọi người lại đẩy tôi ra rìa như
vậy. Có lẽ là vì ngoại hình lạnh lùng của tôi, hoặc cũng có thể là vì
đôi mắt của tôi luôn nhìn xa xăm, và vì miệng của tôi rất ít khi mở lời,
hoặc là vì tôi thích việc quan sát hơn là lao vào chơi đùa.

Và cũng chính vì thế nên tôi đã học cách thu mình lại, xây dựng một bức
tường chia cắt tuyệt đối với thế giới xung quanh, nơi mà lời nói của
người khác chỉ như gió thoảng mây bay. Giao tiếp, đối với tôi nó thật
phiền phức vì nó đòi hỏi tôi vừa phải lộ ra cảm xúc của mình, vừa phải
đoán được cảm xúc của người khác, và tôi không muốn làm vậy một chút
nào. Chính vì thế nên tôi cảm thấy hài lòng với sự cô đơn của mình, và
cảm thấy rằng như vậy là đã ổn rồi.

Ba mẹ rất yêu thương tôi, thật đấy, họ đúng là có bận rộn với công việc,
đúng là thường xuyên bỏ tôi ở nhà một mình giống như trên phim nhưng mỗi
khi có dịp là luôn về thăm tôi, luôn dắt tôi đi chơi. Tôi yêu họ lắm,
chỉ là, tôi đã không còn muốn giao tiếp với mọi người xung quanh nữa.
Việc nói chuyện với người khác thật sự khiến cho tôi cảm thấy rất mệt
mỏi.

Tôi thích bóng tối. Với người khác thì bóng tối có vẻ đáng sợ, nó đầy
rẫy những bóng ma và những nỗi sợ hãi vô tận. Nhưng đối với tôi, bóng
tối giống như một chiếc chăn bông ấm áp, lặng lẽ ôm lấy tôi và bảo vệ
tôi khỏi thế giới xung quanh, một thế giới ồn ào, hỗn loạn, nơi mọi
người cười nói với nhau và mong đợi tôi phải tham gia. Trong bóng tối,
tôi có thể là chính mình, ôm lấy cuốn sách yêu thích và cái đèn pin trên
tay rồi sau đó thỏa sức để trí tưởng tượng bay bổng mà không bị ai quấy
rầy. Bóng tối không phán xét, không đòi hỏi, nó chỉ nhẹ nhàng ôm ấp lấy
tôi như một người bạn thầm lặng vậy.

Tôi thích trời mưa, vì trời mưa giống như là một người ``đồng phạm'' của
tôi vậy. Khi trời mưa, tôi sẽ không phải ra ngoài, không phải nói chuyện
với bất kỳ ai. Chỉ cần một câu nói đơn giản ``Xin lỗi, hôm nay trời
mưa.'' Là tôi có thể từ chối mọi lời mời một cách lịch sự mà không bị
trách móc. Tiếng mưa rơi lộp độp trên mái nhà tạo nên một bản giao hưởng
riêng bệt, và tôi thường ngồi bên cửa sổ để nhìn những giọt nước lăn dài
trên mặt kính, tưởng tượng chúng như những suy nghĩ chảy mãi không cần
đích đến của bản thân.

Thế nhưng rồi có một người đã chìa tay ra với tôi, chọc thủng lớp bảo vệ
mà tôi dựng nên bấy lâu nay và xâm nhập vào thế giới riêng tư của tôi
một cách tự nhiên nhất. Một cô bé tóc trắng hoạt bát mang tên Hikaru.
Ban đầu, tôi chỉ cảm thấy sự xâm nhập của Hikaru thật phiền phức. Cậu ấy
luôn tự ý đến ngồi cạnh tôi, liên tục nói những chuyện trên trời dưới
đất mà không thèm quan tâm đến việc tôi có thèm nghe hay không.

Và rồi cứ thế, lớp một rồi lớp hai, lớp hai rồi lớp ba, dần dần thì tôi
đã quen với sự hiện diện của Hikaru từ lúc nào mà chẳng hay. Cậu ấy
không phán xét sự trầm lặng của tôi, không quan tâm đến việc tôi có đang
nghe hay không. Hikaru chỉ đơn giản là hiện diện ở đó, lặng lẽ trở thành
một phần trong cuộc sống của tôi, kéo tay tôi tham gia vào những trò
chơi ngốc nghếch chỉ có riêng hai đứa. Và tôi cũng bắt đầu có mong muốn
được đáp lại cô gái nhỏ nhắn hoạt bát ấy, đáp lại nguồn sáng duy nhất mà
tôi không muốn tránh né. Thế là Hikaru đã xâm nhập vào cái thế giới nội
tâm được bảo vệ kiên cố của tôi theo một cách đơn giản đến không ai ngờ.

Tôi ghét việc giao tiếp, nó vẫn vô cùng phiền phức đối với tôi. Thế
nhưng đối với Hikaru lại khác, nói chuyện với cậu ấy rất vui vì tôi có
thể dễ dàng bộc lộ cảm xúc của mình và dễ dàng đoán được là Hikaru đang
nghĩ gì.

Tôi thích bóng tối, vì bóng tối cứ như là một thế giới riêng của tôi và
Hikaru vậy. Một thế giới chẳng ai biết. Những buổi tối ngủ cùng nhau
trong căn phòng tối om, thì thầm kể chuyện, chia sẻ những bí mật nhỏ bé
của nhau mà không ai khác được phép nghe. Trong cái bóng tối dịu dàng
ấy, tôi có thể cảm nhận được hơi ấm từ bàn tay của Hikaru nắm chặt lấy
mình, và trái tim của tôi khẽ rung động theo nhịp đập của người bạn
thân.

Tôi thích ánh sáng, vì khi đó tôi sẽ có thể nhìn rõ được khuôn mặt thân
thương của Hikaru. Ánh sáng chiếu lại đôi mắt lấp lánh tuyệt đẹp của
Hikaru, chiếu lên nụ cười rạng rỡ của cậu ấy, chiếu lên khuôn mặt của
một trong những người mà tôi trân trọng nhất trên đời này.

Tôi thích trời mưa, vì khi trời mưa thì Hikaru sẽ dẫn tôi đi tắm mưa.
Hai thân hình nhỏ nhắn của tôi và Hikaru chạy lon ton dưới những giọt
nước mưa nặng hạt xối xả trút xuống, tuy khiến tôi ướt nhẹp và lạnh cóng
nhưng cũng rất vui.

Tôi thích trời nắng. Vì khi trời nắng, tôi và Hikaru sẽ bắt đầu một
chuyến thám hiểm khắp mọi ngóc ngách của thị trấn nhỏ ở quê của hai đứa.
Đi cùng với Hikaru rất vui, dù sau đó hai đứa bọn tôi sẽ lẽo đẽo về nhà
với cái bụng đói meo.

Tôi yêu Hikaru

Thế nhưng tôi không giỏi nói lời yêu, lần đầu tiên tôi căm ghét chính
bản thân vì không thể nào nói những lời hoa mĩ.

Hay nói đúng hơn, tôi thậm chí còn không biết tình yêu là gì cho đến khi
nhận ra rằng mình không thể tưởng tượng nổi một ngày nào đó sẽ không còn
Hikaru bên cạnh.

Tôi sợ.

Sợ rằng nếu nói ra, mọi thứ sẽ thay đổi.

Sợ rằng Hikaru chỉ xem tôi là bạn thân, là người quan trọng, chứ không
phải là người để yêu.

Sợ rằng nếu tôi mở lời, tôi sẽ đánh mất cả mười hai năm trời chúng tôi
đã cùng nhau xây dựng.

Thế nên tôi im lặng.

Tôi để ba người kia tiếp cận.

Tôi để họ nghĩ rằng tôi sẽ chọn một trong số họ.

Tôi để họ nghĩ rằng tôi sắp đưa ra câu trả lời cho một câu chuyện tình
lãng mạn mà ai cũng mong đợi.

Nhưng thật ra, từ đầu đến cuối, trong lòng tôi chỉ có một cái tên duy
nhất.

Hikaru.

Thế nhưng cho dù sợ, cho dù có căm ghét cái vốn từ hạn hẹp của bản thân
đến mức nào thì tôi cũng không thể nào chạy trốn nữa. Tôi và Hikaru sẽ
học ở hai ngành khác nhau, thế nên đây chính là cơ hội cuối cùng để tôi
có thể thổ lộ cảm xúc của bản thân với người tôi yêu.

Tôi không biết bản thân đã phải lòng Hikaru từ khi nào.

Thế nhưng việc đó giờ đây có lẽ đã chẳng còn quan trọng nữa.

Chỉ cần tôi biết rõ rằng tôi yêu cậu ấy từ tận đáy lòng, như vậy là đủ
rồi.

``Bất kể cậu có đang ở đây với tâm trạng thế nào đi nữa, thì hãy nghe
mình nói nhé Hikaru.'' Tôi hít một hơi thật sâu, thổ lộ một trong những
điều quan trọng nhất với một trong những người mà tôi trân trọng nhất
trên đời, ``Mình yêu cậu, yêu cậu rất nhiều! Hãy cho phép mình là người
sẽ ở bên cậu suốt quãng đời còn lại nhé?''

*

Một giây\ldots{}

Hai giây\ldots{}

Và rồi thoáng cái mà đã mười giây trôi qua\ldots{}

Mei nghe thấy tiếng xào xạc phía sau, một âm thanh nhẹ nhàng nhưng rõ
ràng giữa không gian yên tĩnh dưới gốc cây anh đào. Gió chiều thu lay
động những cánh hoa còn sót lại, tạo nên những tiếng động như lời thì
thầm của thiên nhiên, nhưng tiếng xào xạc ấy khác biệt, nó mang theo sự
hiện diện của con người, gần gũi và quen thuộc. Mei dừng lại, trái tim
cô khẽ rung động, không phải vì sợ hãi mà vì một dự cảm mơ hồ.

Cô từ từ quay mặt lại, mái tóc dài đen nhánh khẽ bay theo chuyển động,
và đập vào mắt cô là dáng vẻ của một cô gái thấp hơn cô một chút, mái
tóc trắng ngần để xõa như một dòng suối tươi mát. Đôi mắt màu xanh ngọc
bích như thể chứa đứng cả một hồ nước trong lành bên trong. Và ánh mắt
lấp lánh sự ấm áp, dịu dàng nhưng kiên định khiến Mei bao lần muốn chìm
đắm mãi không thôi.

``Tệ\ldots{} cậu tệ lắm đấy Mei\ldots'' Hikaru cười nói với Mei, giọng
nói vang lên nhẹ nhàng nhưng run run, như một cơn gió thoảng mang theo
nỗi xúc động không thể che giấu.

Nụ cười ấy có chút khác biệt so với nụ cười quen thuộc của Hikaru, nó
mang một nét gì đó trầm lắng hơn và bị mờ đi bởi những giọt nước mắt lăn
dài trên má, long lanh dưới ánh chiều tà, rơi xuống đất như những viên
ngọc. Nước mắt ấy không phải vì buồn bã thuần túy, mà là sự giải tỏa của
bao năm tháng kìm nén, bao cảm xúc dồn nén trong tim. Hikaru lau vội
khóe mi của mình, nhưng nước mắt vẫn tuôn rơi không ngừng, khiến khuôn
mặt cô ấy trở nên mong manh hơn bao giờ hết.

``Mình đã cố gắng lắm mới phớt lờ được thứ tình cảm này\ldots{} để cậu
có thể được hạnh phúc với người khác\ldots{} vậy mà cậu lại\ldots'' Cô
nghẹn lại, không nói hết câu, chỉ lắc đầu, nước mắt càng rơi nhiều hơn.

Lời nói ấy mang theo chút trách móc yêu thương, chút oán hận giả tạo,
nhưng sâu thẳm là sự thừa nhận, thừa nhận rằng tình yêu ấy chưa bao giờ
thực sự biến mất, chỉ là nó luôn bị chôn tít dưới những tầng sâu nhất
bên trong trái tim của Hikaru.

Mei sau đó bước tới một bước, rồi hai bước, khiến cho khoảng cách giữa
hai người dần thu hẹp, và Mei có thể cảm nhận được hơi thở ấm áp của
Hikaru, mùi hương quen thuộc từ mái tóc trắng, một mùi hương nhẹ nhàng,
như hoa anh đào pha lẫn với sự tươi mới của mùa xuân. Mei không nói gì
ngay lập tức, cô chỉ nhìn sâu vào mắt Hikaru, trái tim mình đập mạnh mẽ
hơn bao giờ hết.

Cô tiến lại gần hơn, tay khẽ run khi chạm vào hai bên má của Hikaru, má
cô ấy mềm mại, ấm áp, vẫn còn ướt át vì nước mắt. Mei nâng niu khuôn mặt
ấy như một báu vật quý giá, ngón tay lướt nhẹ qua da thịt mịn màng, xóa
đi những giọt lệ còn đọng lại.

Và rồi, trong một khoảnh khắc dũng cảm mà Mei chưa từng nghĩ mình có, cô
đặt môi mình lên môi Hikaru. Cảm giác này lạ hơn cô nghĩ, mềm mại, ngọt
ngào, như chạm vào một cánh hoa anh đào còn đọng sương mai.

Đôi môi Hikaru đáp lại một cách ngượng ngùng nhưng chân thành, một nụ
hôn đầu tiên đầy vụng về nhưng tràn ngập tình yêu, khiến thế giới xung
quanh dường như tan biến, chỉ còn lại hơi ấm từ nhau. Mei cảm nhận được
vị mặn của nước mắt hòa quyện, vị ngọt của tình cảm dâng trào, và một sự
dễ chịu lan tỏa khắp cơ thể, như thể mọi lớp băng giá trong tim cô đã
tan chảy hoàn toàn. Mei cảm nhận được từng nhịp tim đập loạn của chính
mình, và cả của Hikaru qua lớp áo mỏng.

Nụ hôn của hai người chỉ kéo dài vỏn vẹn vài giây, nhưng với Mei, nó như
một khoảnh khắc vĩnh cửu, một lời khẳng định thầm lặng cho tất cả những
gì cô đã giữ kín bấy lâu.

Mei bỏ môi mình ra, hơi thở dồn dập, và ngay lập tức ôm lấy cơ thể mảnh
mai của Hikaru vào lòng, siết chặt như sợ cô ấy sẽ biến mất. Tay Mei
vuốt ve lưng Hikaru, cảm nhận sự run rẩy nhẹ nhàng từ cơ thể cô ấy, và
mùi hương từ mái tóc trắng lan tỏa, khiến Mei muốn ôm mãi không buông.

``Vậy thì cứ để nó hiện diện đi, không cần phải phớt lờ nó làm gì.'' Mei
thì thầm bên tai Hikaru, giọng nói lạnh lùng thường ngày giờ đây trở nên
ấm áp, dịu dàng đến lạ, như một dòng suối chảy qua trái tim. ``Vì mình
sẽ mãi đồng hành cùng cậu mà.''

Cô ôm Hikaru chặt hơn, cảm nhận nhịp tim hai người hòa quyện, dưới gốc
anh đào chứng kiến bao kỷ niệm, và giờ đây, chứng kiến khởi đầu của một
hành trình mới. Hikaru không trả lời ngay, cô ấy chỉ ôm chặt lấy Mei và
khóc nức nở, tiếng khóc vang lên giữa không gian yên tĩnh, như một dòng
sông cảm xúc cuối cùng cũng được giải phóng. Những giọt nước mắt thấm
ướt vai áo Mei, và Hikaru siết chặt vòng tay, như thể đang ôm lấy toàn
bộ quá khứ, ôm lấy mười hai năm vật lộn với mối tình đơn phương tưởng
chừng như vô vọng, những đêm thao thức tự nhắc nhở bản thân phải buông
bỏ, những khoảnh khắc chạnh lòng khi thấy Mei dần mở lòng với người
khác.

Tiếng khóc ấy không phải đau khổ, mà là sự giải thoát, là niềm vui dâng
trào khi mọi thứ cuối cùng cũng được đáp lại sau tận mười hai năm vô
vọng. Mei vuốt ve tóc Hikaru, để cô ấy khóc hết, và dưới tán anh đào,
hai cô gái ôm nhau, bắt đầu một chương mới của cuộc đời, không còn là
bạn bè thuở nhỏ, mà là những người yêu nhau thực sự, với tình yêu chân
thành và bền vững.

*

``Xong rồi ha?'' Haruka lên tiếng nói với hai người bên cạnh, giọng điệu
không còn khoa trương như mọi khi.

``Ừm, xong rồi. Cả việc giúp đỡ hai người họ, và cả cuộc thi của chúng
ta nữa.'' Shin đáp lại. ``Và tất nhiên, Hikaru chính là người chiến
thắng.''

``Là do tao có vấn đề, hay là do hai thằng bây cứ cố làm ra vẻ nghiêm
túc thế?'' Akira đứng kế bên lên tiếng xỏ xiên.

Cả ba người họ vốn đã biết về tình cảm của Mei ngày hôm qua và đã quyết
định sẽ không đến và chỉ đứng từ xa để quan sát mà thôi. Vì đơn giản là
họ chắc chắn rằng Hikaru sẽ tới, và cũng chắc chắn rằng Mei sẽ nhận ra
sự hiện diện của cô.

Họ từng là những người lạ mặt, sau khi vào ngôi trường này thì trở thành
những đối thủ cạnh tranh nhau từng thứ một từ điểm số đến cả tình cảm
với Mei. Và hiện tại, họ chỉ đơn giản là ba người bạn đang giúp đỡ hai
cô bạn của mình mà thôi.

Sau đó, cả ba người quay lưng rời đi, để lại không gian riêng tư cho Mei
và Hikaru dưới gốc cây anh đào. Bầu trời Tokyo dần chuyển sang sắc tím,
như thể một dấu hiệu khép lại câu chuyện này.
}